\section{はじめに}
\if0
\UTF{2022} テーマの背景
\UTF{2022} システム開発の目的
\UTF{2022} 採用した技術的工夫
\UTF{2022} 得られた結果とそこから導き出せる知見
\fi
オーケストラや器楽合奏などの本格的な合奏体験は，複数の演奏者，楽器，演奏場所を必要とするため，小規模な学校や福祉施設などでは実施が難しい場合がある．特に，児童・生徒が複数の楽器音を組み合わせた合奏を体験する機会は，教育課程や施設環境によって制限されることがある．中学校以降の音楽活動では合唱が中心となる場面も多く，児童・生徒が器楽合奏に触れる機会は小学校段階に偏りやすいことが指摘されている\cite{hira2018}．また，介護施設や福祉施設においても音楽療法や音楽レクリエーションは普及しつつあるが，実施には音楽療法士や演奏者などの人的リソースが必要となる\cite{carelitz_musictherapy}．そのため，専門的な演奏技術を必要とせず，低コストで合奏体験を提供できるシステムには，教育的および福祉的な価値があると考えられる．

そこで本プロジェクトでは，Arduino UNO R4 WiFi を5台用いた小規模なオーケストラシステムを開発する．本システムでは，1台の Arduino を Master，4台の Arduino を Slave として構成し，Master から各 Slave へ I2C 通信によって演奏情報を送信する．各 Slave は，ピアノ，木琴，フルート，ドラムの各楽器を担当し，PC 上の Processing と連携して発音を行う．演奏曲には「森のくまさん」を用い，複数の楽器音を組み合わせながら輪唱形式で演奏する．これにより，少ない機材構成で複数楽器による合奏を再現することを目的とする．

本システムのコンセプトは，「誰でも指揮者になれる」合奏体験である．演奏前には，Processing 上の GUI から BPM，オクターブ，演奏順番を設定できる．また，演奏中にも BPM を変更できるため，使用者はテンポを変化させながら合奏全体を制御できる．これにより，楽譜の知識や楽器の演奏技術がない使用者であっても，GUI 操作を通じて合奏を指揮する体験が可能となる．さらに，Processing 上では楽器音の生成に加えて，BPM に連動してキャラクターの走行速度が変化するクマのアニメーションを実装する．これにより，聴覚的な合奏体験に視覚的な変化を組み合わせ，子供が直感的に楽しめる音楽玩具としての応用を目指す．

本プロジェクトでは，Master から Slave へ演奏順番，オクターブ，BPM などの設定情報を送信し，各 Slave が担当楽器に応じた発音処理を行う構成を実装した．音色の生成には PC 上の Processing を用い，ピアノ，木琴，フルート，ドラムの音色を組み合わせて演奏を行う．また，演奏中の BPM 変更を許可することで，固定テンポの自動演奏ではなく，使用者の操作に応じてテンポが変化する合奏を実現する．

以上より，本システムは，Arduino 複数台と Processing を組み合わせることで，低コストかつ操作しやすい合奏体験を提供できる可能性を示すものである．特に，演奏前の設定変更と演奏中の BPM 変更を可能にしたことで，使用者が合奏全体の進行に関与できる点に特徴がある．このことから，本システムは，子供向けの音楽玩具，学校教育における合奏体験，介護施設や福祉施設での音楽レクリエーションなどへの応用が期待される．
